% Fragmento público generado desde propietarios íntegros.
% Residencia: 52.1.2.
% Fuente: ../incoming/radion_monografia_20260827/manuscrito/sections/02_genealogia.tex:181-272

\subsubsection{Généalogie des constantes et construction du couple angulaire}

\paragraph{Section coinductive et unicité de la valeur générée de \(\pi\)}

Soit \(x_\infty\) la section globale du système projectif. Les lecteurs
coinductifs produisent
\[
\operatorname{ev}_\pi(x_\infty)=\pih,\qquad
\operatorname{ev}_e(x_\infty)=e_{\HMT},\qquad
\operatorname{ev}_\varphi(x_\infty)=\varphi_{\HMT}.
\]
L'indice HMT ne désigne pas un autre nombre. Il conserve la généalogie de la même
constante. En particulier,
\[
\pih=3.14159265358979323846264338327950288419716939937510\ldots.
\]
Le tour positif est
\begin{equation}
T_+(x_\infty)=2\pih.
\label{eq:vuelta-plus}
\end{equation}
Sa publication angulaire d'une unité est
\begin{equation}
R_{\HMT}^{\circ}=\frac{360^\circ}{T_+}=\frac{180^\circ}{\pih}.
\label{eq:radian-publicado}
\end{equation}
L'équation ne définit pas encore le radion enrichi ; elle définit la carte visible
du radian une fois le tour généré.

\paragraph{Bidirectionnalité des chemins et prolongations conjuguées}

La bidirectionnalité ne s'énonce pas comme un slogan. L'involution concrète
\begin{equation}
J_\alpha(T)=\frac{\alphah^2}{T},\qquad
T_-=J_\alpha(T_+),\qquad T_+T_-=\alphah^2
\label{eq:bidireccional-turn}
\end{equation}
conserve un produit et inverse le feuillet. Expansion et contraction sont deux
prolongations conjuguées du même état. La racine finie d'un jet
quadratique peut approcher \(T_+\), mais ne remplace pas la section
coinductive : la différence certifiée est
\[
T_{P_9}-T_+=-5.1110566290335\ldots\times10^{-25}.
\]
Le contrôle fini est un témoin ; le fondement est la limite.

\paragraph{Dépendance causale entre \(\alpha\), action et publication angulaire}

La clôture dodécaphasique génère \(\alphah\) à partir de l'état signé et de ses
cartes réversibles. Ensuite, le lecteur déterminantal local fixe la décennie
d'action au moyen de
\[
\Sigma_H=\varphi_{\HMT}\alphah^{16},
\qquad \operatorname{ord}_{10}\Sigma_H=-34.
\]
La section de retour produit \(\hret\). Ni \(\alpha\) ni \(\hbar\) ne se
\enquote{mesurent en degrés}. Ils sont publiés ultérieurement au moyen du couple angulaire
\begin{equation}
P(\alpha,\eta_{\mathrm{ret}})=
\left([10^3\alpha]_{360},[2(\eta_{\mathrm{ret}}+\alpha)]_{360}\right)
=(A,C^*).
\label{eq:parangular}
\end{equation}
Dans la branche canonique,
\[
A=7.297352569283800997285105472380663\ldots{}^\circ,
\]
\[
C^*=2.1237383389686457242619513803933607937\ldots{}^\circ.
\]
Alors seulement, la carte analytique
\[
x=\frac{\pi A}{180},\qquad y=\frac{\pi C^*}{180}
\]
ouvre les canaux \(q_\pm=e^{-x\mp y}\).

L’ordre causal complet est
\[
\alpha\to(\eta\parallel-34)\to\hbar
\to P\to\kappa_{360}\to q_\pm
\to\text{réalisations MD}.
\]
La relation \(C^*/2-\alpha\) peut servir de vérification inverse. Elle n'est pas utilisée
pour générer \(\hret\) rétroactivement.

\begin{narrativebox}
La carte angulaire ne transforme pas une constante adimensionnelle ou une action en un
angle par décret. Elle publie deux coordonnées de l'état sur une roue finie.
Le degré est le langage de cette publication ; il n'est pas la substance ontologique de
\(\alpha\) ni de \(\hbar\).
\end{narrativebox}
% Fuente: ../incoming/radion_monografia_20260827/manuscrito/sections/03_cierre_exacto.tex:254-285
\paragraph{Indépendance prospective par rapport à la valeur finale}

Le vérificateur de l'opérateur de profondeur interdit expressément trois entrées :
\[
\frac{180}{\pi},\qquad \epsR,\qquad
\text{le décimal cible}.
\]
Il construit d'abord \(S_E\) et \(S_T\), puis exécute la soustraction APP, et ne
compare la clôture qu'à la fin. Les contrôles donnent
\[
\texttt{target\_epsilon\_used=False},\qquad
\texttt{target\_180\_over\_pi\_used=False}.
\]
De même :
\begin{itemize}
\item un seul feuillet donne un défaut d'ordre \(10^{-5}\) ;
\item remplacer \(\theta_2\) par \(\theta_1\) ou \(\theta_3\) déplace le
résultat d'environ \(\pm\pi/6\) ;
\item éliminer le canal \(e\) laisse un défaut d'ordre \(10^{-3}\) ;
\item la variante historique \(\Delta_4(A,C^*)\) produit un autre nombre ;
\item la queue spectrale postérieure à \(\Rstate\) et l'opérateur harmonique à
quatre termes ne coïncident pas non plus avec \(\epsone\).
\end{itemize}

\begin{statusbox}
\textbf{Statut.} La clôture est une \status{formalisation nouvelle} démontrée
avec une structure de départ explicite et un certificat informatique indépendant de la valeur finale.
L'égalité est exacte à la limite coinductive. Le profil court de 36 chiffres
de \(\alpha\) génère une variante qui ne diverge qu'après de nombreux chiffres ;
il est conservé comme témoin historique, non comme valeur canonique.
\end{statusbox}

